% Figure for Calculus §92.
\begin{tikzpicture}[font=\scriptsize]
\begin{axis}[view={145}{24}, axis lines=center, width=9cm, height=8cm,
    xmin=0, xmax=2.6, ymin=0, ymax=1.9, zmin=0, zmax=4.6,
    xtick=\empty, ytick=\empty, ztick=\empty,
    xlabel={$x$}, ylabel={$y$}, zlabel={$z$},
    axis line style={-stealth, thin, black!70},
    every axis x label/.style={at={(axis cs:2.6,0,0)}, anchor=north east},
    every axis y label/.style={at={(axis cs:0,1.9,0)}, anchor=north west},
    every axis z label/.style={at={(axis cs:0,0,4.6)}, anchor=south},
    clip=false]
  % quarter of the paraboloid z = 4 - x^2 - 2y^2 over the first quadrant
  \addplot3[surf, shader=flat, draw=figB!35, fill=figB!12, opacity=0.85,
      domain=0:1, y domain=0:90, samples=11, samples y=10]
      ({2*x*cos(y)}, {sqrt(2)*x*sin(y)}, {4-4*x^2});
  % trace C1 in the plane y = 1
  \addplot3[figB, thick, domain=0:1.4142, samples=40, samples y=0] (x, 1, {2-x^2});
  % trace C2 in the plane x = 1
  \addplot3[figG, thick, domain=0:1.2247, samples=40, samples y=0] (1, x, {3-2*x^2});
  % tangent lines
  \addplot3[figR, very thick, domain=-0.75:0.5, samples=2, samples y=0] ({1+x}, 1, {1-2*x});
  \addplot3[figO, very thick, domain=-0.3:0.25, samples=2, samples y=0] (1, {1+x}, {1-4*x});
  % the point P and its shadow
  \addplot3[black!50, dashed, thin] coordinates {(1,1,0) (1,1,1)};
  \addplot3[only marks, mark=*, mark size=1.4pt, black] coordinates {(1,1,1)};
  \addplot3[only marks, mark=*, mark size=1pt, black!60] coordinates {(1,1,0)};
  \node[anchor=west] at (axis cs:1,1,1.05) {$P(1,1,1)$};
  \node[anchor=north] at (axis cs:1,1,0) {$(1,1)$};
  \node[figB, anchor=south] at (axis cs:0.2,1,1.98) {$C_1$};
  \node[figG, anchor=south west] at (axis cs:1,0.2,2.95) {$C_2$};
  \node[figR, anchor=north] at (axis cs:1.5,1,0) {$T_1$};
  \node[figO, anchor=west] at (axis cs:1,1.25,0) {$T_2$};
\end{axis}
\end{tikzpicture}
